\defmodule {REDBAS}

This module offers tools to reduce a basis in various ways,
and to find a shortest vector in a lattice
using pre-reductions and a branch-and-bound (BB) procedure.
It also has a procedure to compute a Minkowski-reduced basis
(this is much more expensive than finding a shortest vector).
%%  See \cite{rCOU96x} for more details.
One must first call {\tt InitREDBAS} before calling any other procedure
in this module.

\bigskip\hrule
\code

DEFINITION MODULE REDBAS;

IMPORT BASISLR;
FROM LATBASIS IMPORT Basis;
FROM CONFIG   IMPORT BDim;

TYPE
   VerifyType  = (Verify0, Verify2, Verify3);
\endcode
 \tabb Type of verification performed by the tests:
    uses {\tt REDBAS},  {\tt REDBAS2}, or {\tt REDBAS3}
    for {\tt Verify0}, {\tt Verify2}, and {\tt Verify3},
   respectively.
 \endtabb
\code


VAR
   PreRedDieterSV, PreRedLLLSV, PreRedLLLRM : BOOLEAN;
\endcode
  \tabb  These boolean variables indicate which type of pre-reduction 
   is to be performed for {\tt ShortestVector} (SV) and for 
   {\tt ReductMinkowski} (RM).
   {\tt Dieter} means the pairwise pre-reduction as in the procedure
   {\tt PreRedDieter}.  
   {\tt LLL} means the LLL reduction of Lenstra, Lenstra, and Lov\'asz.

   The variable {\tt PreRedDieterSV} is originally set to {\tt TRUE}
   and the 2 others are originally set to {\tt FALSE}.
   These variables are reset automatically depending on the thresholds 
   {\tt MinkLLL, ShortDiet, ShortLLL} as explained below.
   Note that the LLL reduction is never performed when {\tt SISquares = TRUE}.
   \pierre {Reason?}
 \endtabb
\code

   ShortDiet : LONGINT;
\endcode
 \tabb
   Whenever the number of nodes in the BB tree exceeds the threshold 
   {\tt ShortDiet}, in the {\tt ShortestVector} procedure, {\tt PreRedDieterSV}
   is automatically set to {\tt TRUE} for the next call; 
   otherwise it is set to {\tt FALSE}.
   The default value is 1000.
 \endtabb
\code

   ShortLLL : LONGINT;
\endcode
 \tabb 
   Whenever the number of nodes in the BB tree exceeds the threshold 
   {\tt ShortLLL}, in the {\tt ShortestVector} 
   procedure, {\tt PreRedLLLSV}
   is automatically set to {\tt TRUE} for the next call; 
   otherwise it is set to {\tt FALSE}.
   The default value is 1000.
  \endtabb
\code

   MinkLLL : LONGINT;
\endcode
 \tabb  
   Whenever the number of nodes in the BB tree exceeds {\tt MinkLLL}
   in the {\tt ReductMinkowski} procedure,
   {\tt PreRedLLLRM} is automatically set to {\tt TRUE} for the
   next call; otherwise it is set to {\tt FALSE}.
   The default value is 500\,000.
  \endtabb
\code

   MaxPreRed : LONGINT;
\endcode
 \tabb
   Maximum number of transformations in the procedure {\tt PreRedDieter}.
   After {\tt MaxPreRed} successful transformations have been performed,
   the prereduction is stopped.
   The default value is 1\,000\,000.
  \endtabb
\code

   MaxNodesBB : LONGINT;     
\endcode
  \tabb The maximum number of nodes that we are ready to accept
   in the branch-and-bound (BB) tree when calling {\tt ShortestVector}
   or {\tt ReductMinkowski}.  When this number is exceeded, the procedure
   aborts and returns {\tt FALSE}.
   The default value is 10\,000\,000.
  \endtabb
\code


PROCEDURE InitREDBAS (MaxDim : BDim);
\endcode
  \tab  Initializes the internal variables of this 
   module with a given upper dimension {\tt MaxDim}
   for the bases, for memory allocation purposes.
   This {\tt MaxDim} (saved in the internal variable {\tt CurMaxDim})
   should be at least as large as the dimension of any
   {\tt Basis B} that this module has to work on (i.e., {\tt B\^{}.Dim}
   can never exceed {\tt CurMaxDim}).
   \endtab
\code

PROCEDURE DeleteREDBAS;
\endcode
  \tab Deallocates the memory which was allocated when calling
  {\tt InitREDBAS} for the last time.
% for the vector of {\tt SuperIntegers} 
%  {\tt Sv}, the matrices {\tt VTemp}, and 
%  {\tt WTemp}, and the vector {\tt ZeroVect}.
   Must be called only if {\tt InitREDBAS} has been called previously.
  \endtab
\code

PROCEDURE GetMaxDim (): BDim;
\endcode
  \tab  Returns the current value of {\tt CurMaxDim}.
   \endtab
\code

PROCEDURE PairwiseRedPrimal (B : Basis;  i, d : BDim);
\endcode
  \tab  Performs pairwise reductions.
  This procedure tries to reduce each basis vector with 
  index larger than $d$, and distinct from $i$, by adding to it a 
  multiple of the $i$-th vector.  
  Always use the Euclidean norm.
 \endtab
\code

PROCEDURE PairwiseRedDual (B : Basis;  i : BDim);
\endcode
  \tab  Performs pairwise reductions, trying to reduce every other vector of
  the {\em dual\/} basis by adding multiples of the $i$-th vector.
  That may change the $i$-th vector in the primal basis.
  Each such dual reduction is actually performed only if that does not
  increase the length of vector $i$ in the primal basis. 
  Always use the Euclidean norm.
 \endtab
\code

PROCEDURE PreRedDieter (B : Basis;  d : BDim);
\endcode
  \tab  Performs the reductions of the preceding two procedures 
  using cyclically all values of $i$ (only for $i > d$ in the latter case)
  and stops after either {\tt MaxPreRed} successful transformations have been 
  achieved or no further reduction is possible.
  Always use the Euclidean norm.
  \endtab
\code

PROCEDURE SetScale (B : Basis);
\endcode
\tab  This procedure is used when {\tt SISquares = TRUE}, 
   just before doing the Choleski decomposition
   for the branch-and-bound procedure, in the procedures {\tt ShortestVector}
   and {\tt ReductMinkowski}.
   It rescales the square Euclidean lengths of the basis and $m$-dual 
   basis vectors by factors {\tt EchVV} and {\tt EchWW}, respectively,
   and $m^2$ by the product {\tt EchVV * EchWW}.
   These factors are chosen by the procedure 
   so that the log of the largest squared
   Euclidean length is not too far from 0.
   The rescaled values are placed in {\tt VVLR} and {\tt WWLR}, and the 
   rescaled value of $m^2$ in {\tt mLR2}.
   The Choleski decomposition and the branch-and-bound procedures
   use these rescaled values to avoid {\tt LONGREAL} overflow.
\endtab
\code

PROCEDURE RedLLL (B : Basis;  fact : LONGREAL;  m : LONGINT;  Max : BDim);
\endcode
  \tab  Performs a LLL basis reduction with coefficient {\tt fact}, which must
   be smaller than one.  If {\tt fact} is closer to one, the basis will be
   (typically) ``more reduced'', but that will require (slightly) more work.
   Always use the Euclidean norm.
   {\em Danger: With the current implementation, there could be {\tt LONGREAL}
   overflow when {\tt SISquares = TRUE} and there are very long basis vectors.}
  \endtab
\code

PROCEDURE CalculCholeski (B : Basis;  VAR DC2 : BASISLR.BVect;
                          VAR C0 : BASISLR.BMatr) : BOOLEAN;
\endcode
\tab  Performs a Choleski decomposition of the matrix of scalar
  products of the basis vectors, i.e., $\bU'\bU = \bV'\bV$ where $\bV$
  is the matrix of the basis vectors and $\bU$ is upper triangular.
  Returns in {\tt C0} the elements of $\bU$ that are {\em strictly above\/}
  the diagonal, i.e., $u_{i,j}$ for $j > i$.
  Returns in {\tt DC2} the squared elements of the diagonal of $\bU$.
  If {\tt SISquares = TRUE}, these elements are rescaled (divided) by
  {\tt EchVV}.  Returns {\tt TRUE} if the decomposition was successful.
\endtab
\code

PROCEDURE TransformStage3 (B : Basis;  VAR z : BASISLR.VectLI;  VAR k : BDim);
\endcode
 \tab Procedure used in {\tt ReductMinkowski} to perform a transformation
  of stage 3 described in \cite{rAFF85a}. Also used in {\tt ShortestVector}.
  Assumes that $\sum_{i=1}^t z_i V_i$ is a short vector that will enter
  the basis.  Tries to reduce some vectors by looking for indices $i < j$
  such that $|z_j| > 1$ and $q=\lfloor z_i/z_j\rfloor \not=0$,
  and adding $q V_i$ to $V_j$ when this happens.
  Returns in $k$ the last index $j$ such that $|z_j|=1$.
 \endtab
\code

PROCEDURE ShortestVector (B : Basis) : BOOLEAN;
\endcode
  \tab Computes the shortest vector in lattice with basis {\tt B},
   using branch-and-bound.
   If {\tt MaxNodesBB} is exceeded during {\em one} of the branch-and-bounds,
   the procedure aborts and returns {\tt FALSE}.
   Otherwise, it returns {\tt TRUE}.
%  This procedure does not care about numerical imprecision, and therefore,
%  the results are not 100\% reliable.
\endtab
\code

PROCEDURE ReductMinkowski (B : Basis;  d : BDim) : BOOLEAN;
\endcode
   \tab  Reduces {\tt B} into a Minkowski reduced basis, with respect to the
   Euclidean norm, assuming
   that the first $d$ vectors are already reduced and sorted.
   If {\tt MaxNodesBB} is exceeded during {\em one} of the branch-and-bounds,
   the procedure aborts and returns {\tt FALSE}.
   Otherwise, it returns {\tt TRUE}, and the basis is reduced and
   is sorted by vector lengths (the shortest vector is {\tt V[1]} and
   the longest is {\tt V[Dim]}).
   This procedure does not care about numerical imprecision due to the
   (64-bit) floating-point representation.
   In this sense, the results are not 100\% reliable.
   \endtab
\code

END REDBAS.
\endcode


